% GENERATED FILE. Edit canonical lessons, then run npm run generate:books.
% canonical-source-hash: fnv1a64:8eaec7bf7e698c71
% canonical-lessons: GU-C351-hakikat, GU-C351-bhram, GU-C351-gersamaj, GU-C351-yatharth, GU-C351-chetravu

\chapter{Facts and Mix-ups}
\label{ch:gu-facts-and-mix-ups}
\hlchaptermodality{351}

\begin{chapteropening}
\textbf{By the end of this chapter:} \emph{I can tell a fact from a misunderstanding.}

Complete the last lesson of chapter 351: I can tell a fact from a misunderstanding.
\end{chapteropening}

\section[\texorpdfstring{\gu{હકીકત} (hakīkat)}{હકીકત (hakīkat)}]{\gu{હકીકત} (hakīkat) — a fact}

\label{lesson:GU-C351-hakikat}



\begin{quote}
Before the new one: say the Gujarati for seized, then the Gujarati for to recover a debt.
\end{quote}

\subsection*{You'll want to know: \gu{હકીકત}}
\textbf{\gu{હકીકત}} — \emph{hakīkat} — \textquotedblleft{}a fact\textquotedblright{}.

\begin{quote}
\gu{આ} \gu{વાર્તા} \gu{નથી}, \gu{હકીકત} \gu{છે}.
\end{quote}

\begin{quote}
\emph{ā vārtā nathī, hakīkat che.}
\end{quote}

\begin{quote}
This isn't a story; it's a fact.
\end{quote}

A feminine noun.

\begin{cousinweb}
\textbf{\gu{હકીકત}} (\emph{hakīkat}) is usually traced to Arabic \emph{ḥaqīqa}, \textquotedblleft{}truth, reality\textquotedblright{}, through Persian.
\end{cousinweb}

\begin{grammarlens}[title={What you've built}]
One more word for giving reasons and opinions.
\end{grammarlens}

\subsection*{Your turn}
\begin{itemize}
\raggedright
  \item \emph{Say it:} \emph{hakīkat}
  \item \emph{Say it:} \emph{hakīkat}, once more
  \item \emph{Say it:} \emph{ā vārtā nathī, hakīkat che.} — This isn't a story; it's a fact.
  \item \emph{Recall:} say the Gujarati for seized, then the Gujarati for to recover a debt, then say \emph{hakīkat} again
\end{itemize}

\begin{tcolorbox}[breakable,colback=teal!4,colframe=teal!35!black,title={Before you move on}]
What does \gu{હકીકત} mean? (\textquotedblleft{}A fact\textquotedblright{}.) Say it once more.
\end{tcolorbox}
\section[\texorpdfstring{\gu{ભ્રમ} (bhram)}{ભ્રમ (bhram)}]{\gu{ભ્રમ} (bhram) — an illusion}

\label{lesson:GU-C351-bhram}



\begin{quote}
Before the new one: say the Gujarati for to recover a debt, then the Gujarati for a fact.
\end{quote}

\subsection*{You'll want to know: \gu{ભ્રમ}}
\textbf{\gu{ભ્રમ}} — \emph{bhram} — \textquotedblleft{}an illusion\textquotedblright{}.

\begin{quote}
\gu{તે} \gu{સાચું} \gu{નથી}, \gu{તારો} \gu{ભ્રમ} \gu{છે}.
\end{quote}

\begin{quote}
\emph{te sācũ nathī, tāro bhram che.}
\end{quote}

\begin{quote}
It isn't true; you are imagining it.
\end{quote}

A masculine noun.

\begin{cousinweb}
\textbf{\gu{ભ્રમ}} (\emph{bhram}) is a tatsama, borrowed from Sanskrit \emph{bhrama}, \textquotedblleft{}wandering about; confusion\textquotedblright{} (Monier-Williams).
\end{cousinweb}

\begin{grammarlens}[title={What you've built}]
One more word for giving reasons and opinions.
\end{grammarlens}

\subsection*{Your turn}
\begin{itemize}
\raggedright
  \item \emph{Say it:} \emph{bhram}
  \item \emph{Say it:} \emph{bhram}, once more
  \item \emph{Say it:} \emph{te sācũ nathī, tāro bhram che.} — It isn't true; you are imagining it.
  \item \emph{Recall:} say the Gujarati for to recover a debt, then the Gujarati for a fact, then say \emph{bhram} again
\end{itemize}

\begin{tcolorbox}[breakable,colback=teal!4,colframe=teal!35!black,title={Before you move on}]
What does \gu{ભ્રમ} mean? (\textquotedblleft{}An illusion\textquotedblright{}.) Say it once more.
\end{tcolorbox}
\section[\texorpdfstring{\gu{ગેરસમજ} (gersamaj)}{ગેરસમજ (gersamaj)}]{\gu{ગેરસમજ} (gersamaj) — a misunderstanding}

\label{lesson:GU-C351-gersamaj}



\begin{quote}
Before the new one: say the Gujarati for a fact, then the Gujarati for an illusion.
\end{quote}

\subsection*{You'll want to know: \gu{ગેરસમજ}}
\textbf{\gu{ગેરસમજ}} — \emph{gersamaj} — \textquotedblleft{}a misunderstanding\textquotedblright{}.

\begin{quote}
\gu{આ} \gu{નાની} \gu{ગેરસમજને} \gu{કારણે} \gu{બે} \gu{મિત્રો} \gu{લડ્યા}.
\end{quote}

\begin{quote}
\emph{ā nānī gersamajne kāraṇe be mitro laḍyā.}
\end{quote}

\begin{quote}
The two friends fought because of this small misunderstanding.
\end{quote}

A feminine noun.

\begin{cousinweb}
\textbf{\gu{ગેરસમજ}} (\emph{gersamaj}) is built from \emph{ger-}, \textquotedblleft{}mis-, not\textquotedblright{}, usually traced to Persian \emph{ghair}, \textquotedblleft{}other\textquotedblright{}, and \emph{samaj}, \textquotedblleft{}understanding\textquotedblright{}, the noun of \textbf{\gu{સમજવું}} (\emph{samajvũ}), from Sanskrit \emph{saṃbudhyatē} (CDIAL 12959).
\end{cousinweb}

\begin{grammarlens}[title={What you've built}]
One more word for giving reasons and opinions.
\end{grammarlens}

\subsection*{Your turn}
\begin{itemize}
\raggedright
  \item \emph{Say it:} \emph{gersamaj}
  \item \emph{Say it:} \emph{gersamaj}, once more
  \item \emph{Say it:} \emph{ā nānī gersamajne kāraṇe be mitro laḍyā.} — The two friends fought because of this small misunderstanding.
  \item \emph{Recall:} say the Gujarati for a fact, then the Gujarati for an illusion, then say \emph{gersamaj} again
\end{itemize}

\begin{tcolorbox}[breakable,colback=teal!4,colframe=teal!35!black,title={Before you move on}]
What does \gu{ગેરસમજ} mean? (\textquotedblleft{}A misunderstanding\textquotedblright{}.) Say it once more.
\end{tcolorbox}
\section[\texorpdfstring{\gu{યથાર્થ} (yathārth)}{યથાર્થ (yathārth)}]{\gu{યથાર્થ} (yathārth) — true to the facts}

\label{lesson:GU-C351-yatharth}



\begin{quote}
Before the new one: say the Gujarati for an illusion, then the Gujarati for a misunderstanding.
\end{quote}

\subsection*{You'll want to know: \gu{યથાર્થ}}
\textbf{\gu{યથાર્થ}} — \emph{yathārth} — \textquotedblleft{}true to the facts\textquotedblright{}.

\begin{quote}
\gu{આ} \gu{ફોટો} \gu{ગામની} \gu{યથાર્થ} \gu{પરિસ્થિતિ} \gu{બતાવે} \gu{છે}.
\end{quote}

\begin{quote}
\emph{ā phoṭo gāmnī yathārth paristhiti batāve che.}
\end{quote}

\begin{quote}
This photo shows the true state of the village.
\end{quote}

\begin{cousinweb}
\textbf{\gu{યથાર્થ}} (\emph{yathārth}) is a tatsama, borrowed from Sanskrit \emph{yathārtha}, \textquotedblleft{}in accordance with reality\textquotedblright{} (Monier-Williams).
\end{cousinweb}

\begin{grammarlens}[title={What you've built}]
One more word for giving reasons and opinions.
\end{grammarlens}

\subsection*{Your turn}
\begin{itemize}
\raggedright
  \item \emph{Say it:} \emph{yathārth}
  \item \emph{Say it:} \emph{yathārth}, once more
  \item \emph{Say it:} \emph{ā phoṭo gāmnī yathārth paristhiti batāve che.} — This photo shows the true state of the village.
  \item \emph{Recall:} say the Gujarati for an illusion, then the Gujarati for a misunderstanding, then say \emph{yathārth} again
\end{itemize}

\begin{tcolorbox}[breakable,colback=teal!4,colframe=teal!35!black,title={Before you move on}]
What does \gu{યથાર્થ} mean? (\textquotedblleft{}True to the facts\textquotedblright{}.) Say it once more.
\end{tcolorbox}
\section[\texorpdfstring{\gu{છેતરાવું} (chetrāvũ)}{છેતરાવું (chetrāvũ)}]{\gu{છેતરાવું} (chetrāvũ) — to be taken in}

\label{lesson:GU-C351-chetravu}



\begin{quote}
Before the new one: say the Gujarati for a misunderstanding, then the Gujarati for true to the facts.
\end{quote}

\subsection*{You'll want to know: \gu{છેતરાવું}}
\textbf{\gu{છેતરાવું}} — \emph{chetrāvũ} — \textquotedblleft{}to be taken in\textquotedblright{}.

\begin{quote}
\gu{ઘણા} \gu{ગ્રાહકો} \gu{નકલી} \gu{માલથી} \gu{છેતરાય} \gu{છે}.
\end{quote}

\begin{quote}
\emph{ghaṇā grāhko naklī mālthī chetrāy che.}
\end{quote}

\begin{quote}
Many customers are taken in by fake goods.
\end{quote}

Past (he …): \gu{તે} \gu{વેપારીથી} \gu{છેતરાયો}, \emph{te vepārīthī chetrāyo}.

\begin{cousinweb}
\textbf{\gu{છેતરાવું}} (\emph{chetrāvũ}) is the \emph{-āvũ} partner of \textbf{\gu{છેતરવું}} (\emph{chetarvũ}), \textquotedblleft{}to deceive\textquotedblright{}, whose own history is not settled.
\end{cousinweb}

\begin{grammarlens}[title={What you've built}]
One more word for giving reasons and opinions.
\end{grammarlens}

\subsection*{Your turn}
\begin{itemize}
\raggedright
  \item \emph{Say it:} \emph{chetrāvũ}
  \item \emph{Say it:} \emph{chetrāvũ}, once more
  \item \emph{Say it:} \emph{ghaṇā grāhko naklī mālthī chetrāy che.} — Many customers are taken in by fake goods.
  \item \emph{Recall:} say the Gujarati for a misunderstanding, then the Gujarati for true to the facts, then say \emph{chetrāvũ} again
\end{itemize}

\begin{tcolorbox}[breakable,colback=teal!4,colframe=teal!35!black,title={Before you move on}]
What does \gu{છેતરાવું} mean? (\textquotedblleft{}To be taken in\textquotedblright{}.) Say it once more.
\end{tcolorbox}
